% Curs 11 - Pagina 7 (13-XII-1994).
% Fragment pentru % Curs 11 - Pagina 7 (13-XII-1994).
% Fragment pentru % Curs 11 - Pagina 7 (13-XII-1994).
% Fragment pentru % Curs 11 - Pagina 7 (13-XII-1994).
% Fragment pentru \input{courses/pages/curs11-pagina7-ro}.
% Transcriere dupa manuscris. Nu contine preambul sau \newpage.

\subsection*{Curs 11 - Pagina 7}

\noindent\underline{Lem\u{a}:} Fie $m,n,q\in\mathbb N^*$, $n=mq\Rightarrow n\mathbb Z\subseteq m\mathbb Z$.

\[
\frac{m\mathbb Z}{n\mathbb Z}\cong\mathbb Z_q.
\]

\noindent Definim $f$ prin compunerea
\[
\begin{tikzcd}[column sep=large]
m\mathbb Z \arrow[r] \arrow[rr,bend left=30,"f"] &
\mathbb Z \arrow[r,"\mathrm{can}"] &
\mathbb Z/q\mathbb Z=\mathbb Z_q .
\end{tikzcd}
\]
\[
mk\longmapsto k\longmapsto \widehat{k}=k+q\mathbb Z.
\]
\[
m(h+k)\longmapsto h+k\longmapsto\widehat h+\widehat k.
\]

\noindent $f$ e morfism surjectiv $\Rightarrow \operatorname{Im}f=\mathbb Z_q$.
\[
\mathbb Z_q\cong \frac{m\mathbb Z}{\ker f}.
\]

\noindent $x\in\ker f\Rightarrow f(x)=\widehat 0$.
\[
x\in m\mathbb Z,\quad x=mk\ \stackrel{f}{\longmapsto}\ \widehat k=\widehat 0
\iff q\mid k\iff k=qk_1.
\]
\[
\Rightarrow x=mqk_1=nk_1\quad\Rightarrow\quad\ker f=n\mathbb Z.
\]

\bigskip
\begin{center}
\fbox{\textbf{TABLA DE MATERII}}\qquad 2 teorii + 5 exerci\c{t}ii
\end{center}
% Ultima sintagma pare "2 teorii + 5 exercitii"; scrisul este mic.
% Numerotarea manuscrisului are o corectura in dreptul punctului 5:
% atat "legea de compozitie indusa", cat si "monoid, grup" apar cu (5).
% Reproducem dublura pentru a nu renumerota tacit lista originala.

\begin{enumerate}
\item[(1)] Compunerea aplica\c{t}iilor de mul\c{t}imi. Propriet\u{a}\c{t}i.
\item[(2)] Rela\c{t}ii de echivalen\c{t}\u{a}, clase, transversale, mul\c{t}imi factor.
\item[(3)] Teoreme de factorizare pentru aplica\c{t}iile de mul\c{t}imi.
\item[(4)] Legi de compozi\c{t}ie interne; asociativitate, comutativitate, element neutru.
\item[(5)] Legea de compozi\c{t}ie indus\u{a} pe o parte stabil\u{a} \c{s}i pe o mul\c{t}ime factor.
\item[(5)] Monoid, grup, subgrup: defini\c{t}ii, exemple.
\item[(6)] Congruen\c{t}e pe grup. Teorema lui Lagrange.
\item[(7)] Subgrup normal, grup factor.
\item[(8)] Ordinul unui element, propriet\u{a}\c{t}i, exemple.
\item[(9)] Morfisme de grup, exemple, propriet\u{a}\c{t}i.
\item[(10)] Teorema de factorizare a morfismelor de grup, T.~F-I.
% "T. F-I" e prescurtarea din manuscris; probabil teorema fundamentala de izomorfism.
\item[(11)] Grupul $S_n$, transpozi\c{t}ii, cicluri; descomp. unei permut\u{a}ri
\^{i}n produs de transpozi\c{t}ii (cicluri).
% Punctul (11) continua la inceputul paginii 8.
\end{enumerate}
.
% Transcriere dupa manuscris. Nu contine preambul sau \newpage.

\subsection*{Curs 11 - Pagina 7}

\noindent\underline{Lem\u{a}:} Fie $m,n,q\in\mathbb N^*$, $n=mq\Rightarrow n\mathbb Z\subseteq m\mathbb Z$.

\[
\frac{m\mathbb Z}{n\mathbb Z}\cong\mathbb Z_q.
\]

\noindent Definim $f$ prin compunerea
\[
\begin{tikzcd}[column sep=large]
m\mathbb Z \arrow[r] \arrow[rr,bend left=30,"f"] &
\mathbb Z \arrow[r,"\mathrm{can}"] &
\mathbb Z/q\mathbb Z=\mathbb Z_q .
\end{tikzcd}
\]
\[
mk\longmapsto k\longmapsto \widehat{k}=k+q\mathbb Z.
\]
\[
m(h+k)\longmapsto h+k\longmapsto\widehat h+\widehat k.
\]

\noindent $f$ e morfism surjectiv $\Rightarrow \operatorname{Im}f=\mathbb Z_q$.
\[
\mathbb Z_q\cong \frac{m\mathbb Z}{\ker f}.
\]

\noindent $x\in\ker f\Rightarrow f(x)=\widehat 0$.
\[
x\in m\mathbb Z,\quad x=mk\ \stackrel{f}{\longmapsto}\ \widehat k=\widehat 0
\iff q\mid k\iff k=qk_1.
\]
\[
\Rightarrow x=mqk_1=nk_1\quad\Rightarrow\quad\ker f=n\mathbb Z.
\]

\bigskip
\begin{center}
\fbox{\textbf{TABLA DE MATERII}}\qquad 2 teorii + 5 exerci\c{t}ii
\end{center}
% Ultima sintagma pare "2 teorii + 5 exercitii"; scrisul este mic.
% Numerotarea manuscrisului are o corectura in dreptul punctului 5:
% atat "legea de compozitie indusa", cat si "monoid, grup" apar cu (5).
% Reproducem dublura pentru a nu renumerota tacit lista originala.

\begin{enumerate}
\item[(1)] Compunerea aplica\c{t}iilor de mul\c{t}imi. Propriet\u{a}\c{t}i.
\item[(2)] Rela\c{t}ii de echivalen\c{t}\u{a}, clase, transversale, mul\c{t}imi factor.
\item[(3)] Teoreme de factorizare pentru aplica\c{t}iile de mul\c{t}imi.
\item[(4)] Legi de compozi\c{t}ie interne; asociativitate, comutativitate, element neutru.
\item[(5)] Legea de compozi\c{t}ie indus\u{a} pe o parte stabil\u{a} \c{s}i pe o mul\c{t}ime factor.
\item[(5)] Monoid, grup, subgrup: defini\c{t}ii, exemple.
\item[(6)] Congruen\c{t}e pe grup. Teorema lui Lagrange.
\item[(7)] Subgrup normal, grup factor.
\item[(8)] Ordinul unui element, propriet\u{a}\c{t}i, exemple.
\item[(9)] Morfisme de grup, exemple, propriet\u{a}\c{t}i.
\item[(10)] Teorema de factorizare a morfismelor de grup, T.~F-I.
% "T. F-I" e prescurtarea din manuscris; probabil teorema fundamentala de izomorfism.
\item[(11)] Grupul $S_n$, transpozi\c{t}ii, cicluri; descomp. unei permut\u{a}ri
\^{i}n produs de transpozi\c{t}ii (cicluri).
% Punctul (11) continua la inceputul paginii 8.
\end{enumerate}
.
% Transcriere dupa manuscris. Nu contine preambul sau \newpage.

\subsection*{Curs 11 - Pagina 7}

\noindent\underline{Lem\u{a}:} Fie $m,n,q\in\mathbb N^*$, $n=mq\Rightarrow n\mathbb Z\subseteq m\mathbb Z$.

\[
\frac{m\mathbb Z}{n\mathbb Z}\cong\mathbb Z_q.
\]

\noindent Definim $f$ prin compunerea
\[
\begin{tikzcd}[column sep=large]
m\mathbb Z \arrow[r] \arrow[rr,bend left=30,"f"] &
\mathbb Z \arrow[r,"\mathrm{can}"] &
\mathbb Z/q\mathbb Z=\mathbb Z_q .
\end{tikzcd}
\]
\[
mk\longmapsto k\longmapsto \widehat{k}=k+q\mathbb Z.
\]
\[
m(h+k)\longmapsto h+k\longmapsto\widehat h+\widehat k.
\]

\noindent $f$ e morfism surjectiv $\Rightarrow \operatorname{Im}f=\mathbb Z_q$.
\[
\mathbb Z_q\cong \frac{m\mathbb Z}{\ker f}.
\]

\noindent $x\in\ker f\Rightarrow f(x)=\widehat 0$.
\[
x\in m\mathbb Z,\quad x=mk\ \stackrel{f}{\longmapsto}\ \widehat k=\widehat 0
\iff q\mid k\iff k=qk_1.
\]
\[
\Rightarrow x=mqk_1=nk_1\quad\Rightarrow\quad\ker f=n\mathbb Z.
\]

\bigskip
\begin{center}
\fbox{\textbf{TABLA DE MATERII}}\qquad 2 teorii + 5 exerci\c{t}ii
\end{center}
% Ultima sintagma pare "2 teorii + 5 exercitii"; scrisul este mic.
% Numerotarea manuscrisului are o corectura in dreptul punctului 5:
% atat "legea de compozitie indusa", cat si "monoid, grup" apar cu (5).
% Reproducem dublura pentru a nu renumerota tacit lista originala.

\begin{enumerate}
\item[(1)] Compunerea aplica\c{t}iilor de mul\c{t}imi. Propriet\u{a}\c{t}i.
\item[(2)] Rela\c{t}ii de echivalen\c{t}\u{a}, clase, transversale, mul\c{t}imi factor.
\item[(3)] Teoreme de factorizare pentru aplica\c{t}iile de mul\c{t}imi.
\item[(4)] Legi de compozi\c{t}ie interne; asociativitate, comutativitate, element neutru.
\item[(5)] Legea de compozi\c{t}ie indus\u{a} pe o parte stabil\u{a} \c{s}i pe o mul\c{t}ime factor.
\item[(5)] Monoid, grup, subgrup: defini\c{t}ii, exemple.
\item[(6)] Congruen\c{t}e pe grup. Teorema lui Lagrange.
\item[(7)] Subgrup normal, grup factor.
\item[(8)] Ordinul unui element, propriet\u{a}\c{t}i, exemple.
\item[(9)] Morfisme de grup, exemple, propriet\u{a}\c{t}i.
\item[(10)] Teorema de factorizare a morfismelor de grup, T.~F-I.
% "T. F-I" e prescurtarea din manuscris; probabil teorema fundamentala de izomorfism.
\item[(11)] Grupul $S_n$, transpozi\c{t}ii, cicluri; descomp. unei permut\u{a}ri
\^{i}n produs de transpozi\c{t}ii (cicluri).
% Punctul (11) continua la inceputul paginii 8.
\end{enumerate}
.
% Transcriere dupa manuscris. Nu contine preambul sau \newpage.

\subsection*{Curs 11 - Pagina 7}

\noindent\underline{Lem\u{a}:} Fie $m,n,q\in\mathbb N^*$, $n=mq\Rightarrow n\mathbb Z\subseteq m\mathbb Z$.

\[
\frac{m\mathbb Z}{n\mathbb Z}\cong\mathbb Z_q.
\]

\noindent Definim $f$ prin compunerea
\[
\begin{tikzcd}[column sep=large]
m\mathbb Z \arrow[r] \arrow[rr,bend left=30,"f"] &
\mathbb Z \arrow[r,"\mathrm{can}"] &
\mathbb Z/q\mathbb Z=\mathbb Z_q .
\end{tikzcd}
\]
\[
mk\longmapsto k\longmapsto \widehat{k}=k+q\mathbb Z.
\]
\[
m(h+k)\longmapsto h+k\longmapsto\widehat h+\widehat k.
\]

\noindent $f$ e morfism surjectiv $\Rightarrow \operatorname{Im}f=\mathbb Z_q$.
\[
\mathbb Z_q\cong \frac{m\mathbb Z}{\ker f}.
\]

\noindent $x\in\ker f\Rightarrow f(x)=\widehat 0$.
\[
x\in m\mathbb Z,\quad x=mk\ \stackrel{f}{\longmapsto}\ \widehat k=\widehat 0
\iff q\mid k\iff k=qk_1.
\]
\[
\Rightarrow x=mqk_1=nk_1\quad\Rightarrow\quad\ker f=n\mathbb Z.
\]

\bigskip
\begin{center}
\fbox{\textbf{TABLA DE MATERII}}\qquad 2 teorii + 5 exerci\c{t}ii
\end{center}
% Ultima sintagma pare "2 teorii + 5 exercitii"; scrisul este mic.
% Numerotarea manuscrisului are o corectura in dreptul punctului 5:
% atat "legea de compozitie indusa", cat si "monoid, grup" apar cu (5).
% Reproducem dublura pentru a nu renumerota tacit lista originala.

\begin{enumerate}
\item[(1)] Compunerea aplica\c{t}iilor de mul\c{t}imi. Propriet\u{a}\c{t}i.
\item[(2)] Rela\c{t}ii de echivalen\c{t}\u{a}, clase, transversale, mul\c{t}imi factor.
\item[(3)] Teoreme de factorizare pentru aplica\c{t}iile de mul\c{t}imi.
\item[(4)] Legi de compozi\c{t}ie interne; asociativitate, comutativitate, element neutru.
\item[(5)] Legea de compozi\c{t}ie indus\u{a} pe o parte stabil\u{a} \c{s}i pe o mul\c{t}ime factor.
\item[(5)] Monoid, grup, subgrup: defini\c{t}ii, exemple.
\item[(6)] Congruen\c{t}e pe grup. Teorema lui Lagrange.
\item[(7)] Subgrup normal, grup factor.
\item[(8)] Ordinul unui element, propriet\u{a}\c{t}i, exemple.
\item[(9)] Morfisme de grup, exemple, propriet\u{a}\c{t}i.
\item[(10)] Teorema de factorizare a morfismelor de grup, T.~F-I.
% "T. F-I" e prescurtarea din manuscris; probabil teorema fundamentala de izomorfism.
\item[(11)] Grupul $S_n$, transpozi\c{t}ii, cicluri; descomp. unei permut\u{a}ri
\^{i}n produs de transpozi\c{t}ii (cicluri).
% Punctul (11) continua la inceputul paginii 8.
\end{enumerate}
